\section{Evaluation Metrics}
\label{sec:metrics}

Every model in this project, from the persistence baseline through
XGBoost, is scored by the same set of metrics, implemented once in
\texttt{src/evaluate.py} and reused unchanged in every later session. This
section states each metric formally.

\subsection{Mean Absolute Error (MAE)}

\begin{equation}
\label{eq:mae}
\mathrm{MAE} = \frac{1}{n}\sum_{i=1}^{n} \left| y_i - \hat y_i \right|
\end{equation}

MAE is the average size of a miss, in the original units (Mt CO$_2$). It is
interpretable, but because errors are measured in absolute units, its
value is dominated by the largest countries by emissions level.

\subsection{Root Mean Squared Error (RMSE)}

\begin{equation}
\label{eq:rmse}
\mathrm{RMSE} = \sqrt{\frac{1}{n}\sum_{i=1}^n \left(y_i - \hat y_i\right)^2}
\end{equation}

Squaring each error before averaging penalizes large misses quadratically.
Consequently $\mathrm{RMSE} \gg \mathrm{MAE}$ for a given model signals
that its error budget is concentrated in a small number of large misses (a
``fat tail'') rather than spread evenly across rows.

\subsection{Median Absolute Error (MedianAE)}

\begin{equation}
\label{eq:medianae}
\mathrm{MedianAE} = \operatorname{median}_i \left| y_i - \hat y_i \right|
\end{equation}

MedianAE is immune to the largest emitters: it reports the typical error
for the typical country. Comparing Equations~\ref{eq:mae}
and~\ref{eq:medianae} on the same model reveals how much of the average
error is attributable to a handful of very large countries rather than to
the rest of the sample, a scale-dependence issue discussed generally by
\cite{hyndman2006another}.

\subsection{Median Absolute Percentage Error (MdAPE)}

\begin{equation}
\label{eq:mdape}
\mathrm{MdAPE} = 100 \times \operatorname{median}_{i \,:\, y_i \ge \tau}
\left( \frac{|y_i - \hat y_i|}{y_i} \right)
\end{equation}

where $\tau$ is a minimum-level threshold (this project uses
$\tau = 1$~Mt). Percentage errors are undefined or explosive as
$y_i \to 0$; the exclusion in Equation~\ref{eq:mdape} is part of the
metric's definition, not an afterthought, following the general caution
around percentage-based accuracy measures raised by
\cite{hyndman2006another}.

\subsection{Persistence Skill Score}

\begin{equation}
\label{eq:skill}
\mathrm{Skill} = 1 - \frac{\mathrm{MAE}_{\mathrm{model}}}{\mathrm{MAE}_{\mathrm{persistence}}}
\end{equation}

The skill score in Equation~\ref{eq:skill} measures a model's error
relative to the persistence baseline (Section~\ref{sec:baselines}), rather
than in absolute terms. $\mathrm{Skill}=0$ means the model is exactly as
accurate as assuming no change; $\mathrm{Skill}=1$ is the unreachable
limit of zero error; negative skill means the model is less accurate than
doing nothing and must be reported as a loss, not merely a smaller win.

\subsection{The same-rows rule}

Every comparison table in this project scores all candidate models on an
identical set of rows: a row missing any model's prediction, or missing
the target itself, is dropped for every model before
Equations~\ref{eq:mae}--\ref{eq:skill} are computed
(\texttt{evaluate.comparison\_table}). Without this rule, a model could
appear to win only because it was scored on an easier subset of rows.
